\documentclass[12pt]{article}
\usepackage{amsmath, amssymb}
\usepackage[margin=1in]{geometry}
\setlength{\parskip}{6pt}
\title{QUBO Formulation: Nurse Scheduling Assignment Problem}
\date{}
\begin{document}
\maketitle

\subsection*{Nurse Scheduling Assignment Problem}

\paragraph{Problem Statement.}
Given a set of $N$ nurses and $S$ shifts, along with a cost matrix $C = (c_{i,j})$ where $c_{i,j} \ge 0$ denotes the cost of assigning nurse $i$ to shift $j$, the goal is to assign exactly one nurse to each shift while ensuring that each nurse works at most one shift. The objective is to minimize the total assignment cost.

\paragraph{Decision Variables.}
For each nurse $i \in \{0, \dots, N-1\}$ and each shift $j \in \{0, \dots, S-1\}$, introduce a binary decision variable $x_{i,j} \in \{0, 1\}$. The variable $x_{i,j} = 1$ if nurse $i$ is assigned to shift $j$, and $x_{i,j} = 0$ otherwise.

\paragraph{QUBO Derivation.}
\textbf{Step 1 --- Original Objective.}
The problem seeks to minimize the total assignment cost:
\begin{align}
\min_{x} \quad f(x) = \sum_{i=0}^{N-1} \sum_{j=0}^{S-1} c_{i,j} x_{i,j}.
\end{align}

\textbf{Step 2 --- Constraint Formulation.}
The problem is subject to two sets of constraints:
\begin{align}
\sum_{i=0}^{N-1} x_{i,j} &= 1, \quad \forall j \in \{0, \dots, S-1\} \quad \text{(shift coverage)}, \\
\sum_{j=0}^{S-1} x_{i,j} &\le 1, \quad \forall i \in \{0, \dots, N-1\} \quad \text{(nurse capacity)}.
\end{align}

\textbf{Step 3 --- Penalty Term Expansion.}
We convert the constraints into quadratic penalty terms using weights $\lambda_{\text{eq}} > 0$ and $\lambda_{\text{leq}} > 0$.

For the equality constraint, the penalty for shift $j$ is:
\begin{align}
P_j^{\text{eq}}(x) = \lambda_{\text{eq}} \left( \sum_{i=0}^{N-1} x_{i,j} - 1 \right)^2.
\end{align}
Expanding using the idempotent property $x_{i,j}^2 = x_{i,j}$:
\begin{align}
\left( \sum_{i=0}^{N-1} x_{i,j} - 1 \right)^2 &= \sum_{i=0}^{N-1} x_{i,j}^2 - 2 \sum_{i=0}^{N-1} x_{i,j} + 1 + \sum_{\substack{i,k=0 \\ i \neq k}}^{N-1} x_{i,j} x_{k,j} \nonumber \\
&= \sum_{i=0}^{N-1} x_{i,j} - 2 \sum_{i=0}^{N-1} x_{i,j} + 1 + 2 \sum_{0 \le i < k \le N-1} x_{i,j} x_{k,j} \nonumber \\
&= 1 - \sum_{i=0}^{N-1} x_{i,j} + 2 \sum_{0 \le i < k \le N-1} x_{i,j} x_{k,j}.
\end{align}
Thus, the total equality penalty contribution is:
\begin{align}
P^{\text{eq}}(x) = \lambda_{\text{eq}} \sum_{j=0}^{S-1} \left( 1 - \sum_{i=0}^{N-1} x_{i,j} + 2 \sum_{0 \le i < k \le N-1} x_{i,j} x_{k,j} \right).
\end{align}

For the inequality constraint, we use the standard quadratic penalty for $\le 1$:
\begin{align}
P_i^{\text{leq}}(x) = \lambda_{\text{leq}} \left( \sum_{j=0}^{S-1} x_{i,j} \right)^2.
\end{align}
Expanding using $x_{i,j}^2 = x_{i,j}$:
\begin{align}
\left( \sum_{j=0}^{S-1} x_{i,j} \right)^2 &= \sum_{j=0}^{S-1} x_{i,j}^2 + 2 \sum_{0 \le j < l \le S-1} x_{i,j} x_{i,l} \nonumber \\
&= \sum_{j=0}^{S-1} x_{i,j} + 2 \sum_{0 \le j < l \le S-1} x_{i,j} x_{i,l}.
\end{align}
Thus, the total inequality penalty contribution is:
\begin{align}
P^{\text{leq}}(x) = \lambda_{\text{leq}} \sum_{i=0}^{N-1} \left( \sum_{j=0}^{S-1} x_{i,j} + 2 \sum_{0 \le j < l \le S-1} x_{i,j} x_{i,l} \right).
\end{align}

\textbf{Step 4 --- Combined QUBO Hamiltonian.}
The full QUBO energy function $H(x)$ is the sum of the objective and penalty terms:
\begin{align}
H(x) = \sum_{i=0}^{N-1} \sum_{j=0}^{S-1} c_{i,j} x_{i,j} + P^{\text{eq}}(x) + P^{\text{leq}}(x).
\end{align}

\textbf{Step 5 --- Q Matrix Entries and Offset.}
Mapping $x_{i,j}$ to a single index $k = iS + j$ for $k \in \{0, \dots, NS-1\}$, we collect coefficients for linear and quadratic terms. The QUBO takes the form $H(x) = \sum_{k=0}^{NS-1} \sum_{l=0}^{NS-1} Q_{k,l} x_k x_l + c$. The non-zero entries are given by:
\begin{align}
Q_{k,k} &= c_{i,j} - \lambda_{\text{eq}} + \lambda_{\text{leq}}, \quad \text{where } k = iS + j, \\
Q_{k,k'} &= 2\lambda_{\text{eq}}, \quad \text{where } k = iS + j, \ k' = i'S + j, \ i \neq i', \\
Q_{k,k'} &= 2\lambda_{\text{leq}}, \quad \text{where } k = iS + j, \ k' = iS + l, \ j \neq l, \\
c &= S \lambda_{\text{eq}}.
\end{align}
All other $Q_{k,l}$ entries are zero.

\paragraph{Penalty Weight Justification.}
Let $C_{\max} = \max_{i,j} |c_{i,j}|$. To ensure that any infeasible assignment has strictly higher energy than the optimal feasible assignment, the penalty weights must dominate the maximum possible objective value. A violation of the shift coverage constraint changes the constraint sum by at least $1$, incurring a penalty of at least $\lambda_{\text{eq}}$. Since the maximum possible reduction in the objective function is bounded by $N \cdot C_{\max}$, choosing $\lambda_{\text{eq}} > N \cdot C_{\max}$ guarantees that equality violations are penalized more than any feasible cost saving. For the capacity constraint, assigning a nurse to two shifts yields a penalty of $2\lambda_{\text{leq}}$ from the quadratic expansion. Choosing $\lambda_{\text{leq}} > 2 \cdot C_{\max}$ ensures that pairwise inequality violations are penalized more than the sum of any two assignment costs. These bounds guarantee that the global minimum of $H(x)$ lies within the feasible region.

\paragraph{Correctness Argument.}
Minimizing $H(x)$ over $\{0,1\}^{NS}$ recovers the optimal nurse-shift schedule because $H(x)$ is constructed as the original objective plus non-negative penalty terms that vanish exactly on the feasible set. For any feasible assignment, $P^{\text{eq}}(x) = P^{\text{leq}}(x) = 0$, so $H(x)$ equals the assignment cost. For any infeasible assignment, at least one penalty term is strictly positive. By the choice of $\lambda_{\text{eq}}$ and $\lambda_{\text{leq}}$, the energy increase from any constraint violation strictly exceeds the maximum possible decrease in the objective function achievable by any feasible assignment. Consequently, the global minimum of $H(x)$ must occur at a feasible point, and among feasible points, $H(x)$ coincides with the original cost function, guaranteeing that the minimizer corresponds to the optimal schedule.

\paragraph{Illustrative Example.}
Consider a small benchmark instance with $N=2$ nurses and $S=2$ shifts. The variables map as $x_{0,0} \to x_0$, $x_{0,1} \to x_1$, $x_{1,0} \to x_2$, $x_{1,1} \to x_3$. Substituting $N=2$ and $S=2$ into the general formulas yields the explicit Hamiltonian:
\begin{align}
H(x) &= (c_{0,0} - \lambda_{\text{eq}} + \lambda_{\text{leq}})x_0 + (c_{0,1} - \lambda_{\text{eq}} + \lambda_{\text{leq}})x_1 \nonumber \\
&\quad + (c_{1,0} - \lambda_{\text{eq}} + \lambda_{\text{leq}})x_2 + (c_{1,1} - \lambda_{\text{eq}} + \lambda_{\text{leq}})x_3 \nonumber \\
&\quad + 2\lambda_{\text{eq}}(x_0 x_2 + x_1 x_3) + 2\lambda_{\text{leq}}(x_0 x_1 + x_2 x_3) + 2\lambda_{\text{eq}}.
\end{align}
The corresponding Q matrix entries and offset are:
\begin{align}
Q_{0,0} &= c_{0,0} - \lambda_{\text{eq}} + \lambda_{\text{leq}}, & Q_{1,1} &= c_{0,1} - \lambda_{\text{eq}} + \lambda_{\text{leq}}, \\
Q_{2,2} &= c_{1,0} - \lambda_{\text{eq}} + \lambda_{\text{leq}}, & Q_{3,3} &= c_{1,1} - \lambda_{\text{eq}} + \lambda_{\text{leq}}, \\
Q_{0,2} &= Q_{2,0} = 2\lambda_{\text{eq}}, & Q_{1,3} &= Q_{3,1} = 2\lambda_{\text{eq}}, \\
Q_{0,1} &= Q_{1,0} = 2\lambda_{\text{leq}}, & Q_{2,3} &= Q_{3,2} = 2\lambda_{\text{leq}}, \\
c &= 2\lambda_{\text{eq}}.
\end{align}
All other $Q_{k,l}$ entries are zero. Minimizing this polynomial over $\{0,1\}^4$ yields the optimal assignment. For instance, if $c_{0,0}=1, c_{0,1}=4, c_{1,0}=3, c_{1,1}=2$ and $\lambda_{\text{eq}}=5, \lambda_{\text{leq}}=3$, the feasible assignments are $(x_0,x_1,x_2,x_3) \in \{(1,0,0,1), (0,1,1,0)\}$. Evaluating $H(x)$ gives energies $1+2+2(5)+2(3)+10 = 26$ and $4+3+2(5)+2(3)+10 = 31$, respectively. The global minimum occurs at $x_0=1, x_1=0, x_2=0, x_3=1$, corresponding to assigning nurse $0$ to shift $0$ and nurse $1$ to shift $1$, with a minimum total cost of $3$.

\end{document}